% Discrete-Time Control Barrier Functions: definition and forward invariance.
% Compile: pdflatex discrete_time_cbf.tex   (CI compiles all files in this dir)
\documentclass[11pt]{article}

\usepackage{amsmath,amssymb,amsthm}
\usepackage[margin=1in]{geometry}
\usepackage{hyperref}

\newtheorem{definition}{Definition}
\newtheorem{theorem}{Theorem}
\newtheorem{proposition}{Proposition}
\newtheorem{lemma}{Lemma}
\newtheorem{remark}{Remark}

\title{Discrete-Time Control Barrier Functions}
\author{Ali-Eimaan}
\date{\today}

\begin{document}
\maketitle

\begin{abstract}
We define the discrete-time control barrier function (DCBF) and prove the forward-invariance
result that underpins every safety guarantee in this repository. Two repository components depend
on the definition: \texttt{codegen/models.py::dcbf\_constraint} implements the linear
specialisation, and \texttt{MpcCbfSolver::barrierValue} evaluates the squared-distance barrier
that is the repository's single definition of $h$.
\end{abstract}

\section{Setup and notation}
Consider the discrete-time system $x_{k+1} = f(x_k, u_k)$ with $f : \mathcal{X} \times
\mathcal{U} \to \mathcal{X}$ locally Lipschitz, state $x_k \in \mathcal{X} \subseteq
\mathbb{R}^n$ and input $u_k \in \mathcal{U} \subseteq \mathbb{R}^m$. The safe set is
$\mathcal{C} = \{x \in \mathcal{X} : h(x) \ge 0\}$ for a continuously differentiable barrier
$h : \mathcal{X} \to \mathbb{R}$. A function $\alpha : \mathbb{R}_{\ge 0} \to \mathbb{R}_{\ge 0}$
is class-$\mathcal{K}$ if it is continuous, strictly increasing, and $\alpha(0) = 0$. All symbols
are identical to \texttt{docs/README\_math.md}.

\section{Definition}
\begin{definition}[Discrete-time CBF]
$h$ is a discrete-time control barrier function (DCBF) for $x_{k+1} = f(x_k, u_k)$ on
$\mathcal{C}$ if there exists a class-$\mathcal{K}$ function $\alpha$ with $\alpha(r) \le r$ for
all $r \ge 0$ such that
\[
\sup_{u \in \mathcal{U}} \left[ h(f(x,u)) - h(x) \right] \ge -\alpha(h(x))
\qquad \text{for all } x \in \mathcal{C}.
\]
\end{definition}

\begin{remark}
The bound $\alpha(r) \le r$ is required, not merely convenient: without it the induction in
Theorem~\ref{thm:invariance} cannot use $h(x_k) \ge 0$ to conclude $h(x_{k+1}) \ge 0$. A
class-$\mathcal{K}$ function with $\alpha(r) > r$ would make the one-step condition weaker than
the invariance claim it is meant to license.
\end{remark}

\section{Forward invariance}
\begin{theorem}[Invariance of $\mathcal{C}$]\label{thm:invariance}
If $h$ is a DCBF and, at every step, $u_k$ satisfies the DCBF inequality, then
$x_0 \in \mathcal{C}$ implies $x_k \in \mathcal{C}$ for all $k \ge 0$.
\end{theorem}
\begin{proof}
By induction on $k$. The base case is $x_0 \in \mathcal{C}$ by hypothesis, i.e. $h(x_0) \ge 0$.
For the step, the DCBF inequality evaluated at $x_k$ gives
\[
h(x_{k+1}) = h(f(x_k, u_k)) \ge h(x_k) - \alpha(h(x_k)),
\]
and since $h(x_k) \ge 0$ and $\alpha(r) \le r$ for all $r \ge 0$,
\[
h(x_{k+1}) \ge h(x_k) - h(x_k) = 0.
\]
Thus $x_{k+1} \in \mathcal{C}$, completing the induction.
\end{proof}

\section{Linear class-$\mathcal{K}$ specialisation}
Take $\alpha(r) = \gamma r$ with $\gamma \in (0,1]$. The DCBF condition becomes
\[
h(x_{k+1}) - h(x_k) \ge -\gamma h(x_k)
\quad\Longleftrightarrow\quad
h(x_{k+1}) \ge (1-\gamma)\, h(x_k),
\]
and iterating gives the exponential lower envelope
\[
h(x_k) \ge (1-\gamma)^k h(x_0).
\]
Operationally $\gamma$ is the permitted per-step fractional loss of barrier value: $\gamma \to 0$
is maximally conservative, $\gamma = 1$ permits reaching the boundary in a single step. The
envelope is plotted against the closed-loop $h(x_k)$ trace in the ACC reproduction notebook
(\texttt{reproduction/zeng\_acc2021/reproduce\_acc2021.ipynb}, \S3).

\begin{proposition}[Necessity of $\gamma \le 1$]
If $\gamma > 1$ then $1 - \gamma < 0$, and the condition $h(x_{k+1}) \ge (1-\gamma)\, h(x_k)$
admits $h(x_{k+1}) < 0$ from any $h(x_k) > 0$. Concretely, with $h(x_k) = 1$ the bound reads
$h(x_{k+1}) \ge 1 - \gamma$, which is satisfied by $h(x_{k+1}) = 1 - \gamma < 0$: a trajectory
that merely meets the bound can leave $\mathcal{C}$ in one step. Hence $\gamma > 1$ fails to
guarantee forward invariance.
\end{proposition}

\section{The barrier used in this repository}
For obstacle $j$ with centre $p_j$ and (already inflated) effective radius
$r_j^{\mathrm{eff}} = r_j + r_{\mathrm{ego}} + \delta$,
\[
h_j(x) = \|p(x) - p_j\|^2 - (r_j^{\mathrm{eff}})^2,
\]
where $p(x) = P x$ is the position selector. The squared form is deliberate: it is smooth at
$p(x) = p_j$ (unlike $\|p - p_j\| - r$), quadratic in $x$, and its gradient
\[
\nabla h_j(x) = 2 P^\top (P x - p_j)
\]
is affine in $x$, so the constraint Jacobian is affine and the Gauss-Newton Hessian is exact for
the barrier row. \texttt{TubeMpcCbfSolver::tighteningFor} uses exactly this gradient.

\section{Relaxed decay rate}
With $\omega_k \ge 0$ a decision variable, the constraint becomes
\[
h(x_{k+1}) - h(x_k) \ge -\omega_k \gamma\, h(x_k).
\]
\begin{proposition}[Safety under relaxation]\label{prop:relaxed}
If $h(x_k) \ge 0$ and $\omega_k \gamma \le 1$, then $h(x_{k+1}) \ge 0$.
\end{proposition}
\begin{proof}
\[
h(x_{k+1}) \ge h(x_k) - \omega_k \gamma\, h(x_k) = (1 - \omega_k \gamma)\, h(x_k) \ge 0,
\]
using $\omega_k \gamma \le 1$ and $h(x_k) \ge 0$ in the final inequality.
\end{proof}
The implementation enforces the hypothesis as a hard configuration check
(\texttt{omega\_max} $\times$ \texttt{gamma} $\le 1$ in \texttt{MpcCbfSolver::initialize()}),
not a soft warning: a relaxation violating it would void the very safety it is meant to retain.

\section{Relation to the continuous-time definition}
The continuous-time CBF condition is $\dot h \ge -\alpha(h)$. Its forward-Euler discretisation
with step $\Delta t$ is
\[
\frac{h(x_{k+1}) - h(x_k)}{\Delta t} \ge -\alpha(h(x_k)),
\]
which is exactly the DCBF condition with right-hand side $-\Delta t\, \alpha(h(x_k))$. The
discretisation gap is the Euler remainder: for $h$ with Lipschitz gradient of constant
$L_{\nabla h}$, the one-step consistency error is bounded by
\[
\left| h(x_{k+1}) - h(x_k) - \Delta t\, \dot h(x_k) \right|
\le \frac{L_{\nabla h}}{2} \|x_{k+1} - x_k\|^2 .
\]
For the discrete condition to imply the continuous one up to this error, $\Delta t$ must be small
enough that the remainder is dominated by $\alpha(h)$ near the boundary --- precisely why
\texttt{dt} is a solver parameter rather than a free knob: enlarging $\Delta t$ silently weakens
the guarantee the barrier is meant to provide.

\begin{thebibliography}{9}

\bibitem{agrawal2017}
A.~Agrawal and K.~Sreenath, ``Discrete Control Barrier Functions for Safety-Critical Control of
Discrete Systems with Application to Bipedal Robot Navigation,'' in \emph{Proc. Robotics: Science
and Systems (RSS)}, 2017.

\bibitem{ames2017}
A.~D. Ames, X.~Xu, J.~W. Grizzle, and P.~Tabuada, ``Control Barrier Function Based Quadratic
Programs for Safety Critical Systems,'' \emph{IEEE Trans. Automatic Control}, vol.~62, no.~8, 2017.

\bibitem{ames2019}
A.~D. Ames, S.~Coogan, M.~Egerstedt, G.~Notomista, K.~Sreenath, and P.~Tabuada, ``Control Barrier
Functions: Theory and Applications,'' in \emph{Proc. European Control Conf. (ECC)}, 2019.

\bibitem{zeng2021acc}
J.~Zeng, B.~Zhang, and K.~Sreenath, ``Safety-Critical Model Predictive Control with Discrete-Time
Control Barrier Function,'' in \emph{Proc. American Control Conf. (ACC)}, 2021.

\bibitem{zeng2021cdc}
J.~Zeng, Z.~Li, and K.~Sreenath, ``Enhancing Feasibility and Safety of Nonlinear Model Predictive
Control with Discrete-Time Control Barrier Functions,'' in \emph{Proc. IEEE Conf. Decision and
Control (CDC)}, 2021.

\end{thebibliography}

\end{document}
